\documentclass[10pt, leqno21]{article}

% --- Packages ---
\usepackage[utf8]{inputenc}
\usepackage{amsmath, amssymb, amsthm}
\usepackage{graphicx}
\usepackage{tikz}
\usepackage{enumitem}
\usepackage{xcolor}
\usepackage{pgfplots}
\usepackage{gensymb}
\pgfplotsset{width=10cm,compat=1.9}

% --- Screen-Friendly Geometry (修长无边框模式) ---
% 保持A4标准高度(297mm)，将纸张宽度裁减至160mm。
% 左右各留10mm的边距，这样正文宽度刚好是140mm（原版无geometry的最佳阅读宽度）。
\usepackage[
    paperheight=297mm, 
    paperwidth=160mm, 
    left=10mm,
    right=10mm,
    top=15mm, 
    bottom=15mm
]{geometry}

% --- Eye-Care Background (护眼米色背景) ---
\usepackage{pagecolor}
\definecolor{sepiabg}{RGB}{244,236,220}
\pagecolor{sepiabg}

% --- Screen-Friendly Hyperlinks ---
\usepackage[
    colorlinks=true,
    linkcolor=blue!70!black,
    urlcolor=blue!70!black,
    citecolor=green!50!black
]{hyperref}

% --- Homework Header Setup ---
\usepackage{fancyhdr}
\pagestyle{fancy}
\fancyhf{}
% 宽度恢复了一些，页眉可以写全名字
\lhead{\textbf{Qinghao Hu}} 
\chead{\textbf{MATH 135}} 
\rhead{\textbf{Assignment 1}} 
\cfoot{\thepage} 
\setlength{\headheight}{15pt}
\renewcommand{\headrulewidth}{0.4pt}
\renewcommand{\footrulewidth}{0.4pt}

% --- Homework Environments ---


\theoremstyle{plain}
\newtheorem{theorem}{Theorem}[section]
\newtheorem{lemma}[theorem]{Lemma}
\newtheorem{proposition}[theorem]{Proposition}
\newtheorem{corollary}[theorem]{Corollary}
\newtheorem{problem}{Problem} 


\theoremstyle{definition}
\newtheorem{definition}[theorem]{Definition}
\newtheorem{example}[theorem]{Example}

\theoremstyle{remark}
\newtheorem*{remark}{Remark}

% --- Cambridge-style semantic (unnumbered) environments ---
\theoremstyle{definition}
\newtheorem*{axiom}{Axiom}
\newtheorem*{assumption}{Assumption}
\newtheorem*{notation}{Notation}

\newenvironment{solution}
  {\begin{proof}[\textbf{Solution}]}
  {\end{proof}}

% --- Custom Commands ---
\newcommand{\R}{\mathbb{R}}
\newcommand{\N}{\mathbb{N}}
\newcommand{\Z}{\mathbb{Z}}
\newcommand{\Q}{\mathbb{Q}}
\newcommand{\C}{\mathbb{C}}
\newcommand{\ds}{\displaystyle}
\newcommand{\blue}[1]{\textcolor{blue}{#1}}
\newcommand{\red}[1]{\textcolor{red}{#1}}

\newcommand{\mypic}[3]{%
  \begin{figure}[h!]
    \centering
    \includegraphics[width=#3\textwidth]{#1}
    \caption{#2}
  \end{figure}
}

% --- Lists Configuration ---
\setlist[itemize]{label=--}
\setlist[itemize,2]{label=$\circ$}

\begin{document}
\subsubsection*{Q20}
\begin{proof}
  Suppose for the contradiction, the curve $x^2 + y^2 + 3 = 0$ has rational points. 

  Then, there exist $x, y \in \Q$ that $x^2 + y^2 = 3$.

  Since $x$ and $y$ are rational, we can write them with a common denominator: 

  \begin{align*}
    x = \frac{a}{c} \\
    y = \frac{b}{c}
  \end{align*}

  where, 

  $$\exists a, b, c \in \Z, c \neq 0, gcd(a, b, c) = 1.$$

  Let us sub x and y into the curve, 
  \begin{align*}
    \left(\frac{a}{c}\right)^2 + \left(\frac{b}{c}\right)^2 &= 3 \\
    \frac{a^2}{c^2} + \frac{b^2}{c^2} &= 3 \\
    a^2 + b^2 &= 3c^2
  \end{align*}
  By lemma, we know under modulo 4, if n is odd, $n^2 \equiv 1\pmod{4}$. If n is even, $n^2 \equiv0\pmod{4}$.

  Case 1: $c$ is odd

  We get:
  \[
    a^2 + b^2 \equiv 3\pmod{4}
  \]

  This is impossible to achieve. 

  Case 2: $c$ is even

  We get:
  \[
    a^2 + b^2 \equiv 0\pmod{4}
  \]

  This is possible to achieve when $a$ is even and $b$ is even. 

  However, since $\gcd(a, b, c) = 1$, a, b, c cannot be all even.

  Therefore, curve $x^2 + y^2 - 3 = 0$ does not have rational points. 

\end{proof}

\subsubsection*{Q21}
\begin{proof}
  Suppose, for contradiction, $x^2 + y^2 - 3^k = 0$ has rational solutions for k an odd, positive integer. 

  We get $k = 2n + 1, \exists n \in \N$, according to the definition of odd. 

  Let us sub $k = 2n + 1$ into $x^2 + y^2 - 3^k = 0$,
  \begin{align*}
    x^2 + y^2 &= 3^{2n + 1} \\
    x^2 + y^2 &= 3 \times 3^{2n} \\
    x^2  + y^2  &= 3 \times 9^n \\
    a^2 + b^2 &= 3 \cdot 9^nc^2\pmod{4}\\
    a^2 + b^2  &= 3c^2\pmod{4}
  \end{align*}
\end{proof}

\subsubsection*{Q24}
\begin{proof}
  Suppose, for contradiction, JJJJ
\end{proof}

\end{document}

